\documentclass[../Main.tex]{subfiles}
\begin{document}
\chapter{2018-2019学年微积分（一）（下）期末考试}

\section{单项选择题(每小题 3 分，共 18 分)}
\begin{enumerate}
    \item 设$y_{1}(x),y_{2}(x),y_{3}(x)$都是微分方程$y''+a(x)y'+b(x)y=x^{2}$的解，则以下函数中也是该微分方程的解的是（\qquad）.

          \begin{tasks}[label=\textbf{\Alph*.}, label-width=1.5em, column-sep=1em](1)
              \task $y_{1}+y_{2}+y_{3}$
              \task $\dfrac32 y_{1}-\dfrac12 y_{2}+y_{3}$
              \task $y_{1}-y_{2}$
              \task $2y_{3}$
          \end{tasks}

    \item 以下函数在原点可微的是（\qquad）.

          \begin{tasks}[label=\textbf{\Alph*.}, label-width=1.5em, column-sep=1em](2)
              \task $z=x^{2}+y^{2}$
              \task $z=\sqrt{x^{2}+y^{2}}$
              \task $z=x-y$
              \task $z=\sqrt{xy}$
          \end{tasks}

    \item 设$F(x,y)=0$是一条平面光滑曲线，则以下说法中正确的是（\qquad）.

          \begin{tasks}[label=\textbf{\Alph*.}, label-width=1.5em, column-sep=1em](2)
              \task $\{F_{x},F_{y}\}$是该曲线的切矢量
              \task $\{F_{y},F_{x}\}$是该曲线的法矢量
              \task $\{-F_{y},F_{x}\}$是该曲线的切矢量
              \task $\{-F_{x},F_{y}\}$是该曲线的法矢量
          \end{tasks}

    \item 设平面区域$D$由$x^{2}+y^{2}\leq1$表示，区域$D_{1}$是$D$在第一象限的部分，则
          \[
              \iint\limits_{D}(xy+\cos x\sin y)\,\mathrm{d}\sigma=（\qquad）.
          \]

          \begin{tasks}[label=\textbf{\Alph*.}, label-width=1.5em, column-sep=1em](2)
              \task $0$
              \task $4\iint\limits_{D_{1}}xy\,\mathrm{d}\sigma$
              \task $4\iint\limits_{D_{1}}\cos x\sin y\,\mathrm{d}\sigma$
              \task $4\iint\limits_{D_{1}}(xy+\cos x\sin y)\,\mathrm{d}\sigma$
          \end{tasks}

    \item \begin{minipage}[t]{0.93\linewidth}
              设区域$\Omega$由$z=x^{2}+y^{2}$与$z=1$围成，\(I=\iiint\limits_{\Omega}f(z)\,\mathrm{d}v\)，则以下表达式错误的是（\qquad）.\par\smallskip
              \begin{tasks}[label=\textbf{\Alph*.}, label-width=1.5em, column-sep=1em](1)
                  \task $\displaystyle I=\pi\int_{0}^{1}zf(z)\,\mathrm{d}z$ \qquad
                  \task $\displaystyle I=2\pi\int_{0}^{1}r\,\mathrm{d}r\int_{r^{2}}^{1}f(z)\,\mathrm{d}z$
                  \task $\displaystyle I=\int_{-1}^{1}\,\mathrm{d}x\int_{-\sqrt{1-x^{2}}}^{\sqrt{1-x^{2}}}\,\mathrm{d}y\int_{x^{2}+y^{2}}^{1}f(z)\,\mathrm{d}z$
                  \task $\displaystyle I=2\pi\int_{0}^{1}f(z)\,\mathrm{d}z\int_{0}^{z}r\,\mathrm{d}r$
              \end{tasks}
          \end{minipage}

    \item 已知级数$\sum\limits_{n=1}^{\infty}a_{n}^{2}$收敛，则以下说法中正确的是（\qquad）.
          \begin{tasks}[label=\textbf{\Alph*.}, label-width=1.5em, column-sep=1em](2)
              \task $\displaystyle\sum_{n=1}^{\infty}a_n$收敛
              \task $\displaystyle\sum_{n=1}^{\infty}a_na_{n+1}$收敛
              \task $\displaystyle\sum_{n=1}^{\infty}|a_n|$收敛
              \task $\displaystyle\sum_{n=1}^{\infty}(-1)^na_n$收敛
          \end{tasks}
\end{enumerate}

\section{填空题(每小题 4 分，共 16 分)}
\begin{enumerate}
    \setcounter{enumi}{6}
    \item 微分方程$y''+2y'+y=x+2$的通解为\underline{\hspace{5em}}.
          \vspace{1em}

    \item 设$z=f(x^{2}+y^{2},xy)$，其中$f$有连续偏导数，则$\dfrac{\partial z}{\partial x}=$\underline{\hspace{5em}}.
          \vspace{1em}

    \item 设$f(x,y,z)=x^{3}+y^{5}+z^{4}$，则$\text{div }\textbf{grad}f=$\underline{\hspace{5em}}.
          \vspace{1em}

    \item $\displaystyle\sum_{n=1}^{\infty}\frac{2^{n}}{n3^{n}}=$\underline{\hspace{5em}}.
          \vspace{1em}
\end{enumerate}

\section{基本计算题(每小题 7 分，共 42 分)}
\begin{enumerate}
    \setcounter{enumi}{10}
    \item 求点$A(1,2,3)$到直线$L:\dfrac{x-6}{-1}=\dfrac{y-1}{1}=\dfrac{z-6}{-2}$的距离.
          \vspace{11em}

    \item 设方程组
          \[
              \begin{cases}
                  x^{2}-2x+y^{2}+\ln z+z^{3}=1, \\
                  x+y+z=1
              \end{cases}
          \]
          在包含点$(0,0,1)$的一个邻域上确定隐函数$y=y(x),z=z(x)$，求\( \left.\frac{\mathrm{d}y}{\mathrm{d}x}\right|_{x=0}\)，\(\left.\frac{\mathrm{d}z}{\mathrm{d}x}\right|_{x=0} \).
          \vspace{8em}

    \item 求函数$f(x,y)=4x+2xy-x^{2}-3y^{2}$的极值.
          \vspace{8em}

    \item 求$I=\iiint\limits_{\Omega}(x+y+z)\,\mathrm{d}x\,\mathrm{d}y\,\mathrm{d}z$，其中$\Omega$是三坐标面与平面$x+y+z=1$所围成的四面体.
          \vspace{9em}

    \item 求曲线积分$I=\int_{L}(2xy-y^{2}\cos x)\,\mathrm{d}x+(x^{2}-2y\sin x)\,\mathrm{d}y$，其中曲线$L$沿抛物线$y=\pi x-x^{2}+1$从$A(0,1)$到$B(\pi,1)$.
          \vspace{9em}

    \item 判断级数$\displaystyle\sum_{n=1}^{\infty}\dfrac{2+(-3)^{n}}{(2n-1)!!}$的敛散性，若收敛指出是绝对收敛还是条件收敛，并说明理由.
          \vspace{9em}
\end{enumerate}

\section{应用题(每小题 7 分，共 14 分)}
\begin{enumerate}
    \setcounter{enumi}{16}
    \item 求曲面$z=x^{2}+y^{2},0\leq z\leq2$的面积.
          \vspace{10em}

    \item 计算曲面积分$I=\iint\limits_{S}3xz\,\mathrm{d}y\,\mathrm{d}z+yz\,\mathrm{d}z\,\mathrm{d}x-z^{2}\,\mathrm{d}x\,\mathrm{d}y$，其中$S$是曲面$z=x^{2}+y^{2}$与$z=2-x^{2}-y^{2}$所围立体的表面外侧.
          \vspace{10em}
\end{enumerate}

\section{综合题(每小题 5 分，共 10 分)}
\begin{enumerate}
    \setcounter{enumi}{18}
    \item 设$L$是圆周曲线$(x-1)^{2}+(y-1)^{2}=1$（取正向），$f(x)$为正值连续函数。证明：
          \[
              \oint\limits_{L}xf(y)\,\mathrm{d}y-f(x)\,\mathrm{d}x\geq2\pi.
          \]
          \vspace{10em}

    \item 证明：当$-\pi<x<\pi$时，$\sum_{n=1}^{\infty}\frac{(-1)^{n-1}\sin nx}{n}=\frac{x}{2}$.
          \vspace{10em}
\end{enumerate}

\chapter{2018-2019学年微积分（一）（下）期末考试参考答案}
\section{单项选择题(每小题 3 分，共 18 分)}
\begin{enumerate}
    \item \textbf{Solution}. B.
    \item \textbf{Solution}. A.
    \item \textbf{Solution}. C.
    \item \textbf{Solution}. A.
    \item \textbf{Solution}. D.
    \item \textbf{Solution}. B.
\end{enumerate}

\section{填空题(每小题 4 分，共 16 分)}
\begin{enumerate}
    \setcounter{enumi}{6}
    \item \textbf{Solution}. $y=C_{1}e^{-x}+C_{2}xe^{-x}+x$.
    \item \textbf{Solution}. $2xf_{1}+yf_{2}$.
    \item \textbf{Solution}. $6x+20y^{3}+12z^{2}$.
    \item \textbf{Solution}. $\ln3$.
\end{enumerate}

\section{基本计算题(每小题 7 分，共 42 分)}
\begin{enumerate}
    \setcounter{enumi}{10}
    \item \textbf{Solution}. 取直线$L$上一点$B(6,1,6)$，则
          \[
              \overrightarrow{AB}=\{5,-1,3\},\qquad \boldsymbol{s}=\{-1,1,-2\}.
          \]
          所求距离
          \[
              d=\frac{|\overrightarrow{AB}\times\boldsymbol{s}|}{|\boldsymbol{s}|}
              =\frac{|\,\{-1,7,4\}\,|}{\sqrt6}= \sqrt{11}.
          \]

    \item \textbf{Solution}. 对方程组对$x$求导，得
          \[
              \begin{cases}
                  2x-2+2yy'+\dfrac{z'}{z}+3z^{2}z'=0, \\
                  1+y'+z'=0.
              \end{cases}
          \]
          在$(0,0,1)$处化为
          \[
              \begin{cases}
                  -2+4z'(0)=0, \\
                  1+y'(0)+z'(0)=0.
              \end{cases}
          \]
          解得
          \[
              y'(0)=-\frac32,\qquad z'(0)=\frac12.
          \]

    \item \textbf{Solution}. 令
          \[
              \begin{cases}
                  f_{x}=4+2y-2x=0, \\
                  f_{y}=2x-6y=0.
              \end{cases}
          \]
          解得唯一驻点$(3,1)$.
          又
          \[
              f_{xx}=-2,\quad f_{yy}=-6,\quad f_{xy}=2,\quad \Delta=f_{xx}f_{yy}-f_{xy}^{2}=8>0,
          \]
          且$f_{xx}<0$，故在$(3,1)$处取得极大值
          \[
              f(3,1)=6.
          \]

    \item \textbf{Solution}. 由轮换对称性，
          \[
              I=3\iiint\limits_{\Omega}z\,\mathrm{d}v.
          \]
          截面$z=\text{常数}$与$\Omega$相交所得三角形面积为$\dfrac{(1-z)^{2}}{2}$，故
          \[
              I=3\int_{0}^{1}z\cdot \frac{(1-z)^{2}}{2}\,\mathrm{d}z=\frac18.
          \]

    \item \textbf{Solution}. 因
          \[
              Q_{x}=2x-2y\cos x,\qquad P_{y}=2x-2y\cos x,
          \]
          所以被积表达式是某个二元函数的全微分.
          直接凑微分可得原函数为$F(x,y)=x^{2}y-y^{2}\sin x$，故
          \[
              I=F(\pi,1)-F(0,1)=\pi^{2}.
          \]

    \item \textbf{Solution}. 对正项级数
          \[
              \sum_{n=1}^{\infty}\frac{2}{(2n-1)!!}
          \]
          与
          \[
              \sum_{n=1}^{\infty}\frac{3^{n}}{(2n-1)!!}
          \]
          分别用比值判别法，都可判定收敛.
          因而
          \[
              \sum_{n=1}^{\infty}\left|\frac{2+(-3)^{n}}{(2n-1)!!}\right|
          \]
          亦收敛，故原级数绝对收敛.
\end{enumerate}

\section{应用题(每小题 7 分，共 14 分)}
\begin{enumerate}
    \setcounter{enumi}{16}
    \item \textbf{Solution}. 曲面面积元
          \[
              \mathrm{d}S=\sqrt{1+4x^{2}+4y^{2}}\,\mathrm{d}x\,\mathrm{d}y.
          \]
          投影区域为$D:x^{2}+y^{2}\le2$，故
          \[
              S=\iint\limits_{D}\sqrt{1+4x^{2}+4y^{2}}\,\mathrm{d}x\,\mathrm{d}y
              =2\pi\int_{0}^{\sqrt2}r\sqrt{1+4r^{2}}\,\mathrm{d}r
              =\frac{13\pi}{3}.
          \]

    \item \textbf{Solution}. 由 Gauss 公式，
          \[
              I=\iiint\limits_{\Omega}2z\,\mathrm{d}v.
          \]
          改用球坐标，可得
          \[
              I=2\int_{0}^{2\pi}\,\mathrm{d}\theta\int_{0}^{\pi/4}\cos\varphi\sin\varphi\,\mathrm{d}\varphi\int_{0}^{\sqrt2}\rho^{3}\,\mathrm{d}\rho=\pi.
          \]
\end{enumerate}

\section{综合题(每小题 5 分，共 10 分)}
\begin{enumerate}
    \setcounter{enumi}{18}
    \item \textbf{Proof}. 由 Green 公式，
          \[
              \oint\limits_{L}xf(y)\,\mathrm{d}y-f(x)\,\mathrm{d}x
              =\iint\limits_{D}\bigl(f(y)+f(x)\bigr)\,\mathrm{d}x\,\mathrm{d}y.
          \]
          由区域$D:(x-1)^{2}+(y-1)^{2}\le1$的轮换对称性，
          \[
              \iint\limits_{D}f(y)\,\mathrm{d}x\,\mathrm{d}y=\iint\limits_{D}f(x)\,\mathrm{d}x\,\mathrm{d}y.
          \]
          又因$f(x)>0$，从而
          \[
              \oint\limits_{L}xf(y)\,\mathrm{d}y-f(x)\,\mathrm{d}x
              =2\iint\limits_{D}f(x)\,\mathrm{d}x\,\mathrm{d}y\geq 2\iint\limits_{D}1\,\mathrm{d}x\,\mathrm{d}y=2\pi.
          \]

    \item \textbf{Proof}. 将函数$f(x)=\dfrac{x}{2}$在$(-\pi,\pi)$上作$2\pi$周期延拓并展开为傅立叶级数.
          由奇偶性知$a_{n}=0$，而
          \[
              b_{n}=\frac{2}{\pi}\int_{0}^{\pi}\frac{x}{2}\sin nx\,\mathrm{d}x=\frac{(-1)^{n-1}}{n}.
          \]
          因此其正弦级数为
          \[
              \sum_{n=1}^{\infty}\frac{(-1)^{n-1}\sin nx}{n},
          \]
          由傅立叶级数收敛定理可得当$-\pi<x<\pi$时，
          \[
              \sum_{n=1}^{\infty}\frac{(-1)^{n-1}\sin nx}{n}=\frac{x}{2}.
          \]
\end{enumerate}
\end{document}
